%-*-latex-*-
%-*-latex-*-
\newcommand\COURSE{ciss245}
\newcommand\ASSESSMENT{q7201}
\newcommand\ASSESSMENTTYPE{Quiz}
\newcommand\POINTS{\textwhite{xxx/xxx}}

\input{myquizpreamble}
\input{yliow}
\input{\COURSE}
%-*-latex-*-

%-*-latex-*-

%-*-latex-*-
\renewcommand\TITLE{\ASSESSMENTTYPE \ \ASSESSMENT}

\renewcommand\EMAIL{}
\renewcommand\AUTHOR{eagraham1@cougars.ccis.edu}

\textwidth=6in
\begin{document}
\topmatterthree

%------------------------------------------------------------------------------
\nextq
Complete the following bubblesort and swap function.
\\
\ANSWER
\begin{answercode}
#include <iostream>

template < typename T >
void swap(T & a, T & b)
{
  T temp = a;
  a = b;
  b = temp;
}

// Sorts values from address start to address end - 1. Note that
// the value at address end is *not* included in the sort. <-- CORRECTION
// Use the above swap function.
template < typename T >
void bubblesort(T * start, T * end)
{
  for(T * i = start; i < end - 1; i++)
  {
    for (T * j = start; j < end - 1 - (i - start); j++)
    {
      if(*j > *(j + 1))
      {
        swap(j, *(j + 1);
      }
    }
  }
}

int main()
{
    int x[] = {5, 3, 1, 2, 4};
    bubblesort< int >(&x[0], &x[4]); // sorts {5, 3, 1, 2} (4 not included)

    double y[] = {7, 5, 3, 1, 2, 4, 6, 8};
    bubblesort< double >(&y[0], &y[6]);

    return 0;
}
\end{answercode}

%------------------------------------------------------------------------------
\nextq
Complete the following \verb!Ptr! class which is similar to our \verb!IntPtr!
class except that \verb!Ptr! is a class template.
The following test code should work with your \verb!Ptr! class.
Of course there should be no memory issues.
\begin{console}[fontsize=\footnotesize]
#include <iostream>
#include "Ptr.h"

int main()
{
    Ptr< int > p(42);
    std::cout << *p << '\n';

    Ptr< int > q(p);
    std::cout << *q << '\n';

    Ptr< int > r(0);
    r = p;
    std::cout << *r << '\n';

    // Similar test cases for Ptr< double >, Ptr< char >, etc. should work.
    return 0;
}
\end{console}
\\
\ANSWER
\begin{answercode}
# file: Ptr.h

#ifndef PTR_H
#define PTR_H

#include <iostream>

template < typename T >
class Ptr
{
public:
    Ptr(const T & val)

    {
      p_ = new T(val)
    }
    ~Ptr()
    {
      delete p_;
    }
    Ptr(const Ptr< T > & p)

    {
      p_ = new T(*other.p_);
    }
    const Ptr< T > & operator=(const Ptr< T > & p)
    {
      if (this != &p)
      {
        delete p_;
        p_ = new T (*p.p_);
      }
      return *this;
    }
    T & operator*()
    {
      return *p_;
    }
private:
    T * p_;
};

// you need something here
template < typename T >
std::ostream & operator<<(std::ostream & cout, const Ptr< T > & p)
{
  return cout << *p;
}

#endif
\end{answercode}

%------------------------------------------------------------------------------
%-*-latex-*-
\newpage
\input{instructions}
\end{document}

